\begin{lemma}
\label{lemma-dim-1-noetherian-reduced-separated-has-ample}
Let $X$ be a Noetherian reduced separated scheme of dimension $1$.
Then $X$ has an ample invertible sheaf.
\end{lemma}

\begin{proof}
Let $Z_i$, $i = 1, \ldots, n$ be the irreducible components of $X$.
We view these as reduced closed subschemes of $X$.
By Lemma \ref{lemma-dim-1-noetherian-integral-separated-has-ample}
there exist ample invertible sheaves $\mathcal{L}_i$ on $Z_i$.
Set $T = \bigcup_{i \not = j} Z_i \cap Z_j$. As $X$ is Noetherian
of dimension $1$, the set $T$ is finite and consists of closed
points of $X$. For each $i$ we may, possibly after replacing
$\mathcal{L}_i$ by a power, choose $s_i \in \Gamma(Z_i, \mathcal{L}_i)$
such that $(Z_i)_{s_i}$ is affine and contains $T \cap Z_i$, see
Properties, Lemma \ref{properties-lemma-ample-finite-set-in-principal-affine}.

\medskip\noindent
By Lemma \ref{lemma-glue-invertible-sheaves} we can find an invertible sheaf
$\mathcal{L}$ on $X$ and $s \in \Gamma(X, \mathcal{L})$
such that $(\mathcal{L}, s)|_{Z_i} = (\mathcal{L}_i, s_i)$.
Observe that $X_s$ contains $T$ and is set theoretically equal
to the affine closed subschemes $(Z_i)_{s_i}$. Thus it is affine by
Limits, Lemma \ref{limits-lemma-affines-glued-in-closed-affine}.
To finish the proof, it suffices to find for every $x \in X$, $x \not \in T$
an integer $m > 0$ and a section $t \in \Gamma(X, \mathcal{L}^{\otimes m})$
such that $X_t$ is affine and $x \in X_t$. Since $x \not \in T$
we see that $x \in Z_i$ for some unique $i$, say $i = 1$.
Let $Z \subset X$ be the reduced closed subscheme whose underlying
topological space is $Z_2 \cup \ldots \cup Z_n$.
Let $\mathcal{I} \subset \mathcal{O}_X$ be the ideal
sheaf of $Z$. Denote that $\mathcal{I}_1 \subset \mathcal{O}_{Z_1}$
the inverse image of this ideal sheaf under the inclusion
morphism $Z_1 \to X$. Observe that
$$
\Gamma(X, \mathcal{I}\mathcal{L}^{\otimes m}) =
\Gamma(Z_1, \mathcal{I}_1 \mathcal{L}_1^{\otimes m})
$$
see Remark \ref{remark-conductor}. Thus it suffices to find $m > 0$
and $t \in \Gamma(Z_1, \mathcal{I}_1 \mathcal{L}_1^{\otimes m})$
with $x \in (Z_1)_t$ affine. Since $\mathcal{L}_1$ is ample
and since $x$ is not in $Z_1 \cap T = V(\mathcal{I}_1)$
we can find a section
$t_1 \in \Gamma(Z_1, \mathcal{I}_1 \mathcal{L}_1^{\otimes m_1})$
with $x \in (Z_1)_{t_1}$, see Properties, Proposition
\ref{properties-proposition-characterize-ample}.
Since $\mathcal{L}_1$ is ample we can find a section
$t_2 \in \Gamma(Z_1, \mathcal{L}_1^{\otimes m_2})$
with $x \in (Z_1)_{t_2}$ and $(Z_1)_{t_2}$ affine, see
Properties, Definition \ref{properties-definition-ample}.
Set $m = m_1 + m_2$ and $t = t_1 t_2$. Then
$t \in \Gamma(Z_1, \mathcal{I}_1 \mathcal{L}_1^{\otimes m})$
with $x \in (Z_1)_t$ by construction and
$(Z_1)_t$ is affine by Properties, Lemma
\ref{properties-lemma-affine-cap-s-open}.
\end{proof}